Для апробации распознания химических формул проведён тест на десяти изображениях с подробной нотацией и десяти изображениях со скелетной. В тесте используются трех моделей: \textit{MolScribe}, \textit{DECIMER} и \textit{Qwen3VL}. Для каждой модели сравнивались результаты распознавания без предобработки и с применением адаптивной предобработки изображения (см.~таб.~\ref{tab:chem_models_compare}).

\begin{table*}[h]
	\centering
	\small
	\setlength{\tabcolsep}{4pt}
	\renewcommand{\arraystretch}{1.15}
	\resizebox{\textwidth}{!}{
		\begin{tabular}{l c c c c c c}
	\hline
	\multirow{2}{*}{Модель} & \multicolumn{3}{c}{Без предобработки} & \multicolumn{3}{c}{С предобработкой} \\
	\cline{2-7}
	& Подробная & Скелетная & Сумма & Подробная & Скелетная & Сумма \\
	\hline
	MolScribe & 8/10 & 10/10 & 18/20 & 5/10 & 9/10 & 14/20 \\
	DECIMER   & 2/10 & 4/10  & 6/20  & 4/10 & 2/10 & 6/20 \\
	Qwen3VL   & 3/10 & 3/10  & 6/20  & 4/10 & 3/10 & 7/20 \\
	\hline
\end{tabular}
	}
	\caption{Результаты распознавания химических структур в подробной и скелетной нотациях без предобработки и с адаптивной предобработкой}
	\label{tab:chem_models_compare}
\end{table*}

Анализ результатов показывает, что влияние адаптивной предобработки на качество распознавания химических структур оказалось неоднозначным и существенно зависит от используемой модели. Наилучший исходный результат без предобработки продемонстрировала модель \textit{MolScribe}, корректно распознавшая 18 из 20 изображений. После применения предобработки ее итоговый результат снизился до 14 из 20. Наиболее заметное ухудшение наблюдается для подробной нотации, где число правильных распознаваний уменьшилось с 8 до 5. Для скелетной нотации снижение оказалось менее выраженным: с 10 до 9. Это позволяет сделать вывод, что для \textit{MolScribe} адаптивная предобработка в рассматриваемой конфигурации скорее ухудшает входные данные, чем улучшает их.

Модель \textit{DECIMER} показала одинаковый суммарный результат до и после предобработки: 6 корректных распознаваний из 20. Однако распределение качества между типами нотации изменилось. Для подробной записи наблюдается улучшение с 2 до 4 успешных распознаваний, тогда как для скелетной записи результат, напротив, снизился с 4 до 2. Следовательно, предобработка не дала общего прироста качества, но изменила характер восприятия изображения моделью. Это может указывать на то, что отдельные преобразования изображения делают более различимыми детали подробной записи, но одновременно ухудшают восприятие упрощенных скелетных структур.

Модель \textit{Qwen3VL} показала небольшой положительный эффект от предобработки. Общее число корректно распознанных изображений увеличилось с 6 до 7 из 20. Улучшение произошло за счет подробной нотации, где результат вырос с 3 до 4, тогда как для скелетной нотации качество осталось неизменным. Таким образом, для \textit{Qwen3VL} адаптивная предобработка может рассматриваться как умеренно полезный этап, однако ее вклад пока остается ограниченным.

Если рассматривать результаты в совокупности по всем моделям, то без предобработки было корректно распознано 30 из 60 изображений, а с предобработкой --- 27 из 60. Иными словами, суммарный результат после применения адаптивной предобработки оказался ниже. При этом для подробной нотации совокупный результат не изменился, тогда как для скелетной нотации наблюдается общее ухудшение. Это означает, что в рассматриваемой серии экспериментов предобработка не обеспечила универсального повышения качества распознавания химических структур.

Полученные результаты согласуются с тем, что химические изображения предъявляют более жесткие требования к сохранению пространственной структуры объектов, чем обычный текст. Если в случае текстового OCR предобработка может повысить контраст и убрать шум без существенной потери смысла, то для химических формул и молекулярных схем даже небольшое искажение толщины линий, формы связей, положения символов и пространственных отношений между элементами может привести к ухудшению распознавания. Особенно чувствительными к таким изменениям оказываются скелетные структуры, где значительная часть информации задается именно геометрией изображения.

Таким образом, проведенный эксперимент показывает, что адаптивная предобработка не может рассматриваться как универсальный способ повышения качества распознавания химических структур. Для модели \textit{MolScribe} она приводит к заметному ухудшению результата, для \textit{DECIMER} не дает суммарного выигрыша, а для \textit{Qwen3VL} обеспечивает лишь небольшое улучшение. Следовательно, эффективность предобработки зависит как от архитектуры модели, так и от типа химической нотации. По совокупности результатов целесообразно сделать вывод, что применение предобработки для изображений химических структур требует дополнительной настройки и должно рассматриваться как специализированный, а не универсальный этап обработки.